\documentclass[aspectratio=169,10pt,english]{beamer}

\input{../../../_preambulo_beamer.tex}
\input{../../../_comandos_beamer.tex}

\selectlanguage{english}

\title{MA1001B - Statistical Modeling for Decision Making}
\subtitle{Lesson 21: Standard Normal Distribution and Z Scores}
\author{}
\institute{Tecnol\'ogico de Monterrey}
\date{}

\begin{document}

\begin{frame}[plain]
  \vspace{1.2cm}
  {\Large\bfseries MA1001B - Statistical Modeling for Decision Making\par}
  \vspace{0.7cm}
  {\large Lesson 21: Standard Normal and $z$ Scores\par}
  \vspace{0.7cm}
  {\color{TecRojo}\rule{\textwidth}{1.2pt}}
  \vspace{0.8cm}

  MA1001B Faculty\\[0.3cm]
  Term\\[0.3cm]
  Tecnol\'ogico de Monterrey
\end{frame}

\begin{frame}{The Standardization Question}
  \begin{alertblock}{Guiding question}
    How can an audit score of 60 on $N(50,8)$ be read on the single
    standard normal ruler $Z\sim N(0,1)$?
  \end{alertblock}

  \vspace{0.6em}
  By the end of this lesson, you can:
  \begin{itemize}
    \item convert a raw score into $z=(x-\mu)/\sigma$;
    \item read $P(Z\le z)$ from the standard curve;
    \item transfer that area back to $P(X\le x)$;
    \item explain why skew breaks a $z$-score comparison.
  \end{itemize}

  \vfill
  \scriptsize All audit scores and delay times used today are synthetic. No real company data are used.
\end{frame}

\begin{frame}{From a Raw Score to $z$}
  \small
  Operating model: $X\sim N(50,8^2)$. The marked score is $x=60$:
  \[
    z=\frac{60-50}{8}=1.25.
  \]
  One standard deviation from the mean is always $z=\pm 1$:
  \[
    z(42)=-1,\qquad z(58)=1.
  \]

  \vspace{0.3em}
  \begin{exampleblock}{Synthetic audit clock}
    $z$ locates the score in $\sigma$ units. It is not itself a probability.
  \end{exampleblock}
\end{frame}

\begin{frame}[fragile]{Step 1: Compute $z$}
  \begin{columns}[T]
    \begin{column}{0.42\textwidth}
      \small
      \begin{lstlisting}
z = (x - mu) / sigma
z_60 = (60 - 50) / 8
      \end{lstlisting}

      \begin{block}{Verified output}
        $\mu=50$, $\sigma=8$\\
        $z(60)=1.250000$\\
        $z(\mu-\sigma)=-1.000000$\\
        $z(\mu+\sigma)=1.000000$
      \end{block}
    \end{column}
    \begin{column}{0.58\textwidth}
      \centering
      \includegraphics[width=\textwidth]{../figures/standard_normal_01_z_scores.png}
      \vspace{0.2em}
      \scriptsize Operating density from Step 1
    \end{column}
  \end{columns}
\end{frame}

\begin{frame}{The Standard Left Tail}
  \small
  After standardization,
  \[
    P(X\le 60)=P(Z\le 1.25).
  \]
  `scipy.stats.norm.cdf(1.25)` returns $0.894350$. The right tail is
  \[
    P(Z>1.25)=1-0.894350=0.105650.
  \]
\end{frame}

\begin{frame}[fragile]{Step 2: $P(Z\le 1.25)$}
  \begin{columns}[T]
    \begin{column}{0.40\textwidth}
      \small
      \begin{lstlisting}
z_model = norm(0, 1)
p_left = z_model.cdf(1.25)
      \end{lstlisting}

      \begin{block}{Verified output}
        $P(Z\le 1.25)=0.894350$\\
        $P(Z>1.25)=0.105650$\\
        $P(X\le 60)=0.894350$
      \end{block}
    \end{column}
    \begin{column}{0.60\textwidth}
      \centering
      \includegraphics[width=\textwidth]{../figures/standard_normal_02_standard_probabilities.png}
      \vspace{0.2em}
      \scriptsize Standard left tail from Step 2
    \end{column}
  \end{columns}
\end{frame}

\begin{frame}{Step 3: Skew Breaks the Percentile}
  \begin{columns}[T]
    \begin{column}{0.42\textwidth}
      \small
      Exponential delays: mean $=$ sd $=20$.
      Support starts at $z=-1$.

      \vspace{0.4em}
      $P(Z\le -1)=0.158655$ on the bell.\\
      $P(X\le 0)=0.000000$ on the delay clock.

      \vspace{0.4em}
      \textbf{Limit:} $z$ is comparable only when shapes are similar.
      Lesson 22 uses $z$ cuts as specification limits.
    \end{column}
    \begin{column}{0.58\textwidth}
      \centering
      \includegraphics[width=\textwidth]{../figures/standard_normal_03_skew_breaks_z.png}
      \vspace{0.2em}
      \scriptsize Standardized bell versus delay clock from Step 3
    \end{column}
  \end{columns}
\end{frame}

\begin{frame}{Concept Check}
  \begin{enumerate}
    \item Compute $z$ for $x=60$ on $N(50,8)$.
    \item Why does $P(X\le 60)$ equal $P(Z\le 1.25)$?
    \item What probability is $P(Z>1.25)$?
    \item Why is $z=-1$ not the 16th percentile on the delay clock?
  \end{enumerate}

  \vspace{0.6em}
  \begin{block}{Connection to Lesson 22}
    The session can continue directly into specification limits, process
    yield, and a mean shift.
  \end{block}

  \vfill
  \scriptsize Source: OpenStax, \textit{Introductory Business Statistics 2e},
  Chapter 6, Section 6.1. CC BY-NC-SA.
\end{frame}

\appendix



\end{document}
